\section{Dataset Construction Process}
\subsection{Dataset Construction Collection}

% The dataset was collected directly in commercial avocado orchards located in Dak Lak, Lam Dong, and Gia Lai provinces in Vietnam from Jannuary to April 2026. To isolate the leaf region from complex natural backgrounds (soil, grass, sky, other leaves), a portable black fabric backdrop was placed behind each leaf during image capture. This approach preserves real-world lighting conditions (sunlight, cloud cover, shadows) while minimizing background clutter, a common practice in field-based plant phenotyping [citation].

% All images were captured using [thiết bị, e.g., a smartphone camera (iPhone 13) held at approximately 25–35 cm from the leaf surface] with autofocus enabled. No artificial lighting was used; illumination varied naturally across sunny, overcast, and shaded conditions to ensure diversity. The black background also provided high contrast for subsequent leaf segmentation and reduced the need for complex pre‑processing.

% A total of 1836 raw images were acquired from avocado leaves exhibiting six distinct health statuses: healthy, leaf burn, mealybug infestation, leaf curl, [and two more classes, e.g., scab and anthracnose]. Each leaf sample was first inspected by an agricultural expert to confirm disease type before photography.

\subsection{Dataset Annotation}

\subsection{Quality Control}

% ===============================================
% Data Preprocessing
% ===============================================
% \subsection{Data Preprocessing}

% To ensure consistent input dimensions and reduce background variability, all raw images underwent a standard preprocessing pipeline before being used for model training and evaluation. The preprocessing steps are summarised as follows:

% \begin{enumerate}
%     \item \textbf{Center cropping:} Each image was cropped to a square by taking the largest possible square from the centre (i.e., using the shorter side as the crop size). This step removes peripheral background areas while preserving the central leaf region.
%     \item \textbf{Resizing:} The cropped square was resized to a fixed resolution of \(224 \times 224\) pixels using Lanczos interpolation (high-quality downsampling).
%     \item \textbf{Format conversion:} All processed images were saved in PNG format to avoid generation loss and stored in a flat directory structure (\texttt{processed/}) with sequential filenames (\texttt{00001.png}, \texttt{00002.png}, \dots).
% \end{enumerate}

% For model training, we applied two sets of transformations:

% \begin{itemize}
%     \item \textbf{Base transforms (used for validation and test):} Images were normalised using ImageNet channel statistics (mean = [0.485, 0.456, 0.406], standard deviation = [0.229, 0.224, 0.225]) to facilitate transfer learning from pretrained models.
%     \item \textbf{Augmented transforms (used for training, when enabled):} In addition to the base normalisation, we applied random horizontal flipping (probability 0.5), random rotation within \(\pm 15^\circ\), and random brightness/contrast adjustments (\(\pm 20\%\) each) to increase data diversity and improve model generalisation. No geometric distortion (e.g., shear, zoom) was employed to preserve the shape of disease lesions.
% \end{itemize}

% \sloppy
% The preprocessing and augmentation choices were kept intentionally simple to focus on the intrinsic difficulty of the dataset and to provide a reproducible baseline for future comparisons. All transformations were implemented in PyTorch using standard \texttt{torchvision} modules.
